\section{Related Work}

\paragraph{Context attribution for text LLMs.}
A recent line of work attributes a generated statement to the context passages that support it.
ContextCite~\cite{cohen-wang2024contextcite} hides random subsets of the context and fits a sparse linear model of the response probability; ARC-JSD~\cite{li2026arcjsd} and AttriBoT~\cite{liu2025attribot} make such perturbation estimates more accurate or cheaper.
MIRAGE~\cite{qi-etal-2024-model} uses gradient saliency to link answer tokens to retrieved documents.
Closer to our signal, AT2~\cite{cohenwang2025at2} learns one weight per attention head from perturbation data, and AttnTrace~\cite{wang2026attntrace} and \citet{ding2025attention} read attention directly for text contexts.
All of these define sources as discrete text units and are evaluated on text.
OmniTrace uses the same kind of signal as the attention-based methods, needs no training, and is defined over units of any modality, including time intervals.
The term \emph{generation-time} is also used for prompted citation~\cite{saxena2025generation, gao2023alce}; we use it for attribution that is read from the model's internals while it decodes.

\paragraph{Attribution for vision--language models.}
vStream~\cite{kang2026vstream} streams attribution to image regions during reasoning but learns a map from attention to perturbation effects. VTrace~\cite{yuan2026vtrace} and ERCR~\cite{liang2025ercr} are post-hoc methods for image patches, and MultAttnAttrib~\cite{tran2026multattnattrib} applies selected attention heads to long documents.
These methods cover images and text. None covers audio or video timelines, and none tests attribution causally on continuous inputs.

\paragraph{Temporal evidence in audio and video.}
Models can be trained or prompted to state the time intervals that support an answer~\cite{lu2025vited, wang2025groundedvideollm, zhu2026timelens2}.
MuRGAt~\cite{wan2026murgat} shows that such stated citations are often hallucinated even when the reasoning is correct, and NExT-GQA~\cite{xiao2024nextgqa} shows that video models often answer correctly without looking at the relevant moment.
Mechanistic studies trace where audio information enters the answer across layers~\cite{cho2026tracingaudio, chen2026causalaudio}, but they do not attribute individual generated claims.
We measure which intervals a model relied on for each sentence and compare them with the intervals it states.

\paragraph{Evaluating attribution.}
Attribution can be scored by agreement with human evidence or by its causal effect on the output~\cite{deyoung2020eraser, atanasova-etal-2020-diagnostic}. The two can disagree: \citet{alhamidi2026carve} find that the sensitivity of video models to removed frames tracks annotated evidence only weakly.
We therefore use causal tests as the main evaluation and report agreement with human evidence separately.
